
\subsubsection{Interpretation of PoC Results}

The proof-of-concept validation successfully demonstrates the four-stage validation pipeline's methodology while producing provisional quantitative results. We validated that the framework functions (gate logic discriminates, statistical computations execute, visualization generates), but numerical claims remain tentative pending real infrastructure validation.

\textbf{High-Confidence Finding:} CodeBLEU demonstrates measurement reliability. The observed CV=1.39\% is consistent with the metric's deterministic design. Real validation would likely yield CV in the 0-2\% range.

\textbf{Medium-Confidence Finding:} The four-stage pipeline's Stage 1 multi-criteria gate successfully filters unreliable proxies. The scoped success criterion prevents all-or-nothing failure while maintaining rigor.

\textbf{Low-Confidence Finding:} Runtime efficiency proxy requires hardware performance counters to achieve CV $\leq 5$\%. The PoC's 6.22\% CV, combined with COFFE literature reporting 2-3\% CV for instruction counts, suggests the marginal failure is a PoC artifact. However, empirical validation is required.

\subsubsection{Theoretical Implications}

\#\#\#\# Compositional Validation Principle

Results provide empirical support for compositional validation: different quality dimensions have distinct measurement reliability profiles.

\textbf{Structural Metrics:} Deterministic computations $\rightarrow$ CV $\approx 0-2$\%. Negligible measurement noise. These proxies are ''free'' for RL optimization.

\textbf{Efficiency Metrics:} Require specialized instrumentation (hardware performance counters). Wall-clock time is noisy. Only instruction-level metrics achieve CV ~2-3\%.

\textbf{Learned Metrics:} Require training data and model training. Measurement reliability depends on model quality.

\textbf{Practical Consequence:} Multi-objective RL practitioners cannot assume all auxiliary rewards are equally reliable. Correlation with human judgment does not guarantee reliability.

\#\#\#\# Measurement Theory Meets ML Evaluation

This work bridges psychometrics with ML evaluation. The three criteria (CV, Cohen's d, Spearman $\rho$) are standard in construct validity testing for psychological instruments. Our contribution: importing reliability standards from social science to code generation evaluation.

Current ML benchmarking focuses on predictive validity and construct validity via correlation. We add reliability testing: is the metric stable enough to optimize?

\subsubsection{Limitations and Scope Boundaries}

\#\#\#\# PoC Synthetic Data

\textbf{Limitation:} All measurements use synthetic data. No real CodeLlama-7B generation, no HumanEval execution, no hardware counters.

\textbf{Why This Matters:} Quantitative thresholds may shift when measured on real data.

\textbf{Why Acceptable:} Phase 2C experiment design explicitly scoped PoC as methodology validation. The two-phase approach is a principled risk-reduction strategy.

\textbf{Mitigation Plan:} Phase 4 full-scale validation is the immediate next step. Download CodeLlama-7B-Instruct, generate 500 solutions on 50 HumanEval problems, run \texttt{perf stat -e instructions}, compute actual metrics. Expected timeline: 2-4 weeks.

\#\#\#\# Partial Proxy Validation

\textbf{Limitation:} Only 1/3 proxies validated. Multi-objective optimization proceeds with execution + CodeBLEU, not the originally envisioned execution + structure + efficiency + style.

\textbf{Why This Matters:} Downstream hypotheses lose statistical power for multi-proxy interaction analysis.

\textbf{Why Acceptable:} The scoped gate design treats partial validation as scientifically valid. Single validated proxy is sufficient to test the core hypothesis.

\textbf{Mitigation Plan:} Re-test runtime proxy with real \texttt{perf} during h-e2 execution. Defer PR-style to future work.

\#\#\#\# Runtime Proxy Failure Attribution

\textbf{Limitation:} We attribute runtime proxy's marginal failure to PoC synthetic noise, but have not empirically confirmed this claim.

\textbf{Why This Matters:} If real \texttt{perf} measurements also yield CV >5\%, the efficiency dimension is unmeasurable.

\textbf{Why Acceptable:} COFFE's findings are from real hardware, and the CPU time equation provides theoretical grounding. The inference is most parsimonious given available evidence.

\textbf{Mitigation Plan:} Ablation study comparing PoC synthetic noise vs real \texttt{perf} measurements.

\#\#\#\# Single Programming Language

\textbf{Limitation:} All validation is Python-specific. Results may not generalize to other languages.

\textbf{Why Acceptable:} Python is the dominant language in ML/AI code generation research. Multi-language validation is a natural extension.

\textbf{Future Work:} Extend h-e1 to HumanEval-X or MBXP. Test whether thresholds hold across languages.

\subsubsection{Positioning the Contribution}

\textbf{Methodological Contribution (HIGH Confidence):} The four-stage validation pipeline's Stage 1 is a functional, reusable framework applicable to any proxy-based optimization domain.

\textbf{Empirical Contribution (MEDIUM Confidence):} CodeBLEU demonstrates measurement reliability suitable for RL optimization. This is the first systematic reliability validation of a code generation structural metric.

\textbf{Practical Contribution (MEDIUM Confidence):} The PoC validation strategy is resource-efficient for academic labs without large GPU budgets.

\textbf{What Remains Claim-Free:} Numerical thresholds are provisional. Runtime proxy validation requires confirmation. Multi-objective RL efficacy is tested in future hypotheses.

\subsubsection{Broader Impact}

\textbf{Reduced Compute Waste:} Early-stage proxy filtering prevents wasted GPU hours on unmeasurable signals. At scale, this represents significant compute savings and carbon footprint reduction.

\textbf{Open Science Practices:} The framework promotes transparency by requiring measurement reliability reporting before optimization claims.

\textbf{Reproducibility:} The two-phase PoC $\rightarrow$ real validation approach enables methodology verification before infrastructure investment.
